\documentclass[5p,authoryear]{elsarticle}
\makeatletter 
\def\ps@pprintTitle{%
 \let\@oddhead\@empty
 \let\@evenhead\@empty
 \let\@evenfoot\@oddfoot} % Supprimer le bas de page ELSEVIER
\makeatother
\usepackage[utf8]{inputenc} % En unicode
\usepackage[T1]{fontenc}
\usepackage[english]{babel}
\usepackage[babel=true]{csquotes} % permet de faire \enquote{a} (« a »)
\usepackage[fleqn]{amsmath} % pour certains signes mathématiques
\usepackage{amsthm} % Pour \begin{gather}
\usepackage{booktabs} % pour \toprule (un style de tableau)
\usepackage{multirow} % Pour colonnes multiples des tableaux
\usepackage{amssymb} % Pour \leqslant (<=, >=)
\usepackage{float}
\usepackage{hyperref} % DOIT ETRE EN DERNIER
\usepackage[english]{cleveref} % permet de faire \cref au lieu de \ref (DOIT ETRE EN DERNIER)

\bibliographystyle{elsarticle-num}
%\usepackage[style=authoryear,backend=biber]{biblatex}
%\addbibresource{bibliographie.bib}

\begin{document}

\begin{frontmatter}

\title{Reframing Nuclear Licensing: The FSAR as a Tool, Not a Barrier}

%% Group authors per affiliation:
\author{Gemma (Ira) Strong}
\address{Idaho State University\\ Nuclear Engineering Department }

\author{NE 5599-003 }
\address{Professor Dr. Richard Schultz}
\address{May 4th, 2025}

\begin{abstract}
There is a growing perception among advanced reactor startups that the U.S. Nuclear Regulatory Commission imposes unnecessary delays in the licensing process. However, this view often overlooks a key factor: inadequate understanding of the licensing framework, particularly the role of the Final Safety Analysis Report. This white paper repositions the FSAR not as a bureaucratic barrier but as a strategic tool for achieving regulatory approval. Through comparisons with the Production Part Approval Process used in manufacturing, this paper underscores the importance of upfront planning, experienced licensing personnel, and rigorous documentation. Case studies—including NuScale, Oklo, and Kairos Power—highlight how FSAR quality and team preparedness directly impact licensing timelines. Additionally, the paper explores evolving regulatory practices in the U.S. and internationally, including efforts by the IAEA and the NRC’s upcoming Part 53 rule. By shifting the industry’s mindset and strengthening FSAR development practices, startups can align innovation with regulatory expectations and accelerate the deployment of Small Modular Reactors and Generation IV technologies.
\end{abstract}

\begin{keyword}
Final Safety Analysis Report (FSAR) \sep Nuclear Regulatory Commission (NRC) \sep Small Modular Reactors (SMRs) \sep Generation IV Reactors \sep Licensing Process\end{keyword}

\end{frontmatter}

%\footnotetext{xxx}

\section{Introduction}\label{introduction}
\emph{Nuclear Power Plant Licensing} in the United States is often perceived as an onerous process, especially for innovators developing Small Modular Reactors (SMRs) and Generation IV designs. A central component of this process is the Final Safety Analysis Report (FSAR) – a massive document that details the design and safety case of a reactor. The FSAR has a reputation for being a barrier due to its scope and complexity. However, this white paper posits that the FSAR should be viewed as a tool for ensuring safety and facilitating communication between reactor designers and regulators, rather than as an impediment to innovation. To support this perspective, we examine the U.S. Nuclear Regulatory Commission (NRC) licensing framework and the role of the FSAR within it, highlighting best practices, international parallels, and real-world case studies of both successful and challenged licensing efforts. By demystifying the FSAR’s purpose and contents, we aim to show that a well-prepared FSAR can streamline the licensing of advanced reactors, helping projects avoid delays and achieve timely approvals.\\

This paper seeks to reframe the FSAR not as a barrier, but as a strategic asset in achieving safe, timely reactor deployment. We examine the U.S. Nuclear Regulatory Commission (NRC) licensing framework and how a well-prepared FSAR can streamline the process, reduce regulatory friction, and improve safety outcomes. Additionally, we explore recent reforms to the NRC's mandatory hearing procedures and highlight opportunities to improve efficiency through the National Environmental Policy Act (NEPA) process - particularly via categorical exclusions for microreactors and graded environmental reviews. Drawing parallels with other industries, we argue for a cultural shift in how licensing documentation is perceived and managed, moving from checkbox compliance to design-integrated quality assurance.\\

By reframing the FSAR's purpose, contextualizing its regulatory role, and showcasing lessons learned from recent case studies, we aim to equip developers and policymakers with a roadmap for licensing success in the era of advanced nuclear. A process first mindset grounded in transparency, completeness, and iteration — can transform the FSAR and other regulatory steps into levers of innovation, rather than inhibitors of it.


\section{The Value of Documentation in Safety-Critical Industries}
\vspace{0.2cm}

In many industries, including automotive, aerospace, and nuclear energy, complex documentation processes are often viewed as bureaucratic obstacles rather than tools for safety, quality, and transparency. One such example is the Production Part Approval Process (PPAP) used in manufacturing. PPAP is intended to ensure that suppliers understand design specifications and produce parts reliably under real-world conditions. Despite its critical role in maintaining supply chain integrity and product safety, PPAP is frequently seen as burdensome or excessive by teams focused solely on speed or cost. This negative perception mirrors how some reactor developers approach the FSAR - as a regulatory hurdle rather than an integrated part of engineering development.\\

The mindset that documentation requirements are merely "paperwork" to satisfy compliance authorities is not unique to nuclear energy. It is a broader cultural trend across industries that prioritize short-term deliverables over long term accountability. However, just as high-reliability industries have increasingly embraced PPAP and similar quality assurance frameworks as enablers of product integrity and customer trust, the nuclear sector must embrace the FSAR process as a proactive design tool.\\

Both PPAP and FSAR serve similar foundational goals: risk mitigation, performance validation, and cross-functional alignment. In the case of FSAR, it is not just a licensing requirement; it is the culmination of system understanding, safety justification, and operational foresight. When teams approach FSAR preparation as an integral part of engineering much like mature manufacturing teams view PPAP as part of their design for manufacture philosophy the quality of the resulting system and the efficiency of regulatory review are dramatically improved.\\

Changing the mindset toward licensing documentation is essential, especially as newer, leaner companies enter the nuclear space with minimal institutional experience. FSAR success, like PPAP success, depends less on checklists and more on cultural buy-in: documentation is not the enemy; it is the medium through which safety is proven, collaboration is structured, and future scalability is enabled.

\section{NRC Licensing Process Overview} \label{NRC Licensing Process Overview}

The NRC’s licensing process for nuclear reactors is thorough and multi-faceted, involving safety, environmental, and financial reviews, as well as public hearings. For traditional large reactors, licensing was historically a two-step process under 10 CFR Part 50: first a Construction Permit (with a Preliminary Safety Analysis Report, PSAR) and later an Operating License (with an updated FSAR). Today, most new power reactors (including SMRs) use the one-step Combined License (COL) process under 10 CFR Part 52, which authorizes construction and operation in a single license. In the COL process, the applicant submits a complete FSAR up front, along with other required documentation, as part of the combined license application. The NRC then conducts extensive technical reviews, asks questions (formal Requests for Additional Information, or RAIs), and ultimately issues a Safety Evaluation Report if the application is found acceptable. A mandatory hearing and final Commission vote are also required before license issuance.

\subsection{Role of the FSAR in Licensing}
\vspace{0.2cm}

The FSAR is arguably the cornerstone of the COL application. NRC regulations stipulate that "the application must contain a final safety analysis report that describes the facility, presents the design bases and the limits on its operation, and presents a safety analysis of the structures, systems, and components of the facility as a whole". In essence, the FSAR is where the applicant demonstrates that the proposed reactor meets all safety requirements. It must contain sufficient detail to allow the NRC to reach a final conclusion on all safety matters before granting a license. The NRC's content requirements for FSARs are codified in 10 CFR 52.79, which enumerates dozens of specific topics and analyses that must be included \citep{CFR52.79}. These range from site characteristics and design descriptions to accident analyses and technical specifications. In NRC's safety review, each aspect of the FSAR is evaluated against regulatory standards and guidance to ensure that the reactor can be constructed and operated without undue risk to public health and safety.

\subsection{Licensing Step in Brief}
\vspace{0.2cm}

In practice, an applicant often engages in pre-application interactions with NRC staff to discuss the novel aspects of their design and the planned content of the FSAR. This early engagement can include white papers, technical reports, and meetings addressing key safety topics \cite{NIA2024}. Once the application (including the FSAR) is submitted, the NRC staff performs an acceptance review to confirm it has enough information to docket (formally accept) the application. After docketing, the detailed technical review begins, where NRC specialists in various disciplines examine the FSAR chapter by chapter. They will issue RAIs for any missing information or clarifications needed. The applicant’s timely and thorough responses to these RAIs are critical for keeping the review on schedule. When the staff’s questions are resolved, they document their findings in a Final Safety Evaluation Report (FSER). In parallel, the NRC reviews the environmental report (for NEPA compliance) and holds public hearings (mandatory hearings and possibly contested hearings if there are intervenors). Only after all these steps – and the NRC is satisfied that the FSAR demonstrates safety – will a license be granted. This rigorous process can span several years for a first-of-a-kind reactor design, but as we will explore, the duration is highly sensitive to the quality of the initial FSAR and the efficiency of interactions between the applicant and NRC.

\subsection{Categorical Exclusions and Streamlined NEPA Reviews for Microreactors}
\vspace{0.2cm}

The licensing process for microreactors and advanced reactor designs such as SMRs can be substantially accelerated by streamlining the environmental review mandated under the National Environmental Policy Act (NEPA). Historically, NEPA reviews have been rigorous, time consuming, and resource intensive, often extending licensing timelines significantly. Recognizing this challenge, the Idaho National Laboratory (INL) recently explored opportunities for categorical exclusions specific to microreactors, emphasizing that microreactors inherently have lower environmental footprints due to their smaller size, modularity, inherent safety features, and reduced waste outputs \cite{inl2025categorical}. Categorical exclusions under NEPA allow certain federal actions—which do not individually or cumulatively have significant environmental impacts—to bypass the preparation of extensive environmental assessments (EAs) or environmental impact statements (EISs). For microreactors, applying categorical exclusions could mean drastically reducing the scope of required environmental documentation, thereby speeding up the overall licensing timeline.\\

The INL report highlights how such exclusions can be justified scientifically and technically due to the limited scope of impact associated with microreactors. It recommends specific regulatory language and guidance to ensure a clear, consistent implementation of these exclusions \cite{inl2025categorical}. Additionally, it underscores the role of transparency and public communication in maintaining trust while pursuing regulatory efficiency. Moreover, recent changes by the NRC have begun reflecting a more streamlined approach towards mandatory hearings associated with the COL process. This initiative complements the concept of categorical exclusions by further trimming unnecessary procedural overhead. The NRC has clarified expectations, tightened timelines, and aimed at better managing the mandatory hearing process, thus reducing administrative burdens without compromising safety or public participation \cite{hogan2023nrc}.\\

Together, these regulatory innovations—categorical exclusions tailored for microreactors and streamlined mandatory hearing processes—have the potential to significantly reduce licensing timelines and costs. Consequently, reactor developers can allocate resources more efficiently towards rigorous FSAR preparations and safety analyses, aligning with the overarching goal of deploying safe, innovative nuclear technologies more promptly and effectively.\\

\subsection{Recent Improvements to the NRC's Mandatory Hearing Process}
\vspace{0.2cm}

Recognizing the need for more efficient regulatory processes, the NRC has recently streamlined its mandatory hearing procedures for COLs. This is particularly significant for applicants developing SMRs and advanced Generation IV reactors, where timely and predictable licensing is crucial. Historically, the mandatory hearing—required under Section 189(a) of the Atomic Energy Act for all new reactor COL applications—has been identified as a potential source of delay due to its comprehensive nature and procedural complexity \cite{hogan2023nrc}. In response, the NRC implemented reforms aimed at enhancing efficiency without compromising safety or public participation.\\

These reforms include clearer guidelines on the scope of the mandatory hearings, tighter timelines for the NRC's issuance of questions and the applicants' responses, and more structured criteria for what constitutes adequate safety documentation \cite{hogan2023nrc}. Additionally, the NRC has adjusted its internal processes to enable earlier preparation and coordination among technical and legal staff, reducing delays typically associated with the hearing preparation phase. The NRC's approach now emphasizes a focused review of safety-significant issues, ensuring the hearing addresses pertinent safety and environmental matters without unnecessary repetition of documentation already evaluated during technical reviews. This refinement not only shortens review times but also increases transparency, allowing stakeholders to clearly understand the NRC’s decision-making process.\\

These recent procedural enhancements are exemplified by cases like NuScale Power's SMR design certification, where the streamlined mandatory hearing contributed to achieving timely regulatory approval \cite{hogan2023nrc}. The improvements reflect an important shift toward regulatory efficiency, potentially making the U.S. more attractive for innovative nuclear projects while maintaining rigorous safety standards.\\

\section{The Final Safety Analysis Report: Purpose and Content}

Far from being a mere bureaucratic formality, the FSAR serves a vital technical purpose in reactor development and licensing. It is the definitive safety case for the plant, documenting how the design meets regulatory requirements and identifying how the reactor's design features ensure safe operation. Preparing an FSAR is an intensive engineering and analysis effort that often proceeds in parallel with the plant design itself. In fact, the discipline of developing an FSAR can be a tool for the reactor designer, forcing a rigorous examination of safety from every angle and thereby improving the design.

\subsection{What the FSAR Includes?} \label{What the FSAR Includes}
\vspace{0.2cm}

By design, the FSAR is comprehensive. NRC guidance (such as Regulatory Guide 1.206 \cite{RG1.206} and the Standard Review Plan NUREG-0800) lays out a standard format for FSARs, typically broken into chapters covering topics such as:
\vspace{0.2cm}

\begin{itemize}
    \item \textbf{General Plant Description:} An overview of the nuclear power plant's design and major features.
    \item \textbf{Site Characteristics:} Data on the reactor site’s geology, seismology, hydrology, meteorology, and nearby facilities, demonstrating the site's suitability and how external hazards are accounted for in the design.
    \item \textbf{Reactor and System Design:} Detailed descriptions of structures, systems and components (SSCs) – including the reactor core, reactor coolant system, control systems, electrical power systems, cooling systems, containment or confinement structures, and auxiliary systems. For each, the FSAR explains design bases (the fundamental criteria and assumptions), functional requirements, and how the design meets safety criteria.
    \item \textbf{Safety Analyses:} This includes analysis of normal operations, transients, and postulated accident scenarios. The FSAR presents results of thermal-hydraulic analyses, reactor kinetics, containment response, and radiological consequence calculations for design-basis accidents. It identifies the engineered safety features (like emergency core cooling, containment systems) and evaluates their performance in accident conditions to show that regulatory dose limits and safety margins are met.
    \item \textbf{Technical Specifications and Operational Programs:} The proposed limits and conditions for operation (tech specs), which ensure the plant operates safely within the analyzed envelope. The FSAR also describes operational programs such as radiation protection, fire protection, cybersecurity, quality assurance, maintenance and testing programs, and emergency preparedness plans.
    \item \textbf{Management and Staffing:} Descriptions of the organizational structure, training and qualifications of personnel, and how human factors are addressed in the design and operation of the plant.
    \item \textbf{Safety Evaluation Conclusions:} A summary demonstrating that the design, as described and analyzed in the FSAR, meets the NRC’s safety requirements and will protect public health and safety and common defense and security.
\end{itemize}

\vspace{0.2cm}
These elements align with international best practices as well. The International Atomic Energy Agency (IAEA) provides guidance on the format and content of safety analysis reports in Safety Guide SSG-61, which mirrors the NRC’s approach \cite{IAEA2021}. For example, IAEA recommends including chapters on general plant description, site evaluation, design of safety systems, accident analysis, operational limits, radiation protection, and more – all of which correspond to sections found in a U.S. FSAR. This common structure ensures that regulators worldwide are looking for the same kind of information to make safety determinations. 

\subsection{Regulatory Requirements and Best Practices} 
\vspace{0.2cm}

U.S. regulations not only require an FSAR, but also specify its content in high-level terms. According to NRC's 10 CFR 52.79(a), an FSAR for a combined license must contain design descriptions and analyses sufficient for the NRC to conclude that the plant can be built and run safely \cite{CFR52.79}. This includes information on site suitability, design of SSCs with safety functions, evaluation of the plant's response to postulated accidents, and how regulatory criteria (e.g. dose limits in 10 CFR 50.34 or 10 CFR Part 100 for radiological consequences) are satisfied. In short, the FSAR must "show its work" for every aspect of nuclear safety. Preparing such a document requires a coordinated effort by a multidisciplinary team of nuclear engineers, reactor physicists, thermal-hydraulic analysts, materials experts, risk analysts, and others. \\

A high-quality FSAR is characterized by clarity, thoroughness, and traceability of information. Best practices for applicants include:

\vspace{0.2cm}
\begin{itemize}
    \item \textbf{Following Established Guidance:} Applicants typically use NRC's Standard Review Plan (SRP) as a checklist to ensure the FSAR covers all points that NRC reviewers will expect to see. This alignment helps prevent omissions that could trigger additional questions. Many applicants also follow industry guidelines (e.g., Nuclear Energy Institute templates) to present information consistently.
    \item \textbf{Early and Frequent Engagement with NRC:} As noted, engaging the regulator early to discuss novel design features or safety approaches can prevent major surprises. If a reactor uses an unconventional technology (say, a non-light-water coolant or passive safety features), an applicant might submit topical reports or white papers ahead of the FSAR to gain NRC feedback on methodologies. This way, by the time the FSAR is submitted, the approaches used (e.g., certain analytical codes or safety evaluation methods) are already generally agreed upon. Such pre-application interactions were crucial in recent advanced reactor projects \cite{NIA2024}.
    \item \textbf{Ensuring Technical Adequacy:} All analytical results in the FSAR - from neutronics to accident dose projections - should be backed by validated methods and sufficient data. If prototype testing or separate-effects experiments are needed to confirm a novel concept, those should be done (or at least well underway) by the time of FSAR submission. The NRC will scrutinize whether design claims (e.g., "this reactor can cool itself without power for 72 hours") are supported by analysis or testing. Any gaps here will lead to RAIs and potential delays.
    \item \textbf{Internal Peer Review and Consistency Checks:} Given the FSAR's size (often on the order of more than 10,000 pages for large projects), consistency is a challenge. Leading applicants run rigorous internal reviews to ensure that, for example, a pump described in Chapter 5 has the same performance characteristics in the accident analysis of Chapter 15. Inconsistencies or errors in the FSAR can erode the NRC's confidence and result in lengthy Q\&A to resolve.
    \item \textbf{Leveraging Experience and Expertise:} Teams that include veterans of past licensing efforts tend to navigate the process more effectively. NRC staff have noted that applications from organizations new to nuclear regulation can suffer from "regulatory capability gaps" - unfamiliarity with what level of detail or rigor is expected \cite{Sam2023}. As a result, those applications might omit important analyses or documentation, prolonging the review. On the other hand, companies that invest in experienced licensing engineers and advisors are often better at producing an FSAR that meets NRC's expectations on the first try.
\end{itemize}

\vspace{0.2cm}
It is important to note that the FSAR is not static. Even after a license is issued, the FSAR (now often called the Updated FSAR) must be maintained throughout the life of the plant to reflect changes, as required by 10 CFR 50.71(e). This ensures that the document remains a living reference for the plant's licensing basis. However, for the purposes of initial licensing of SMRs and Gen IV reactors, our focus is on the FSAR as part of the application.

\section{FSAR Impact on SMRs and Generation IV Reactors}

New reactor concepts - including SMRs (which are often scaled-down light-water reactors or advanced modular designs) and Gen IV reactors (which may use novel coolants like gas, liquid metal, or salt, and advanced fuel cycles) - present both opportunities and challenges in the licensing arena. Many SMR and Gen IV developers are enthusiastic about inherent safety features and simplifications in their designs. But these novel features must still be translated into a convincing FSAR to satisfy regulators. Here we explore how the FSAR process affects these advanced technologies, and how it can be managed to not become a roadblock.


\subsection{Unique Challenges for Advanced Designs} \label{Unique Challenges for Advanced Designs}
\vspace{0.2cm}

Traditional light-water reactors (LWRs) have a licensing framework built over decades of operating experience and research. In contrast, a Generation IV sodium-cooled fast reactor or a high-temperature gas reactor might have little prior licensing precedent. This means the applicant's FSAR needs to address safety questions that NRC reviewers may not have encountered before, potentially requiring more foundational explanation. For example, an SMR that uses passive safety cooling (with no operator action or AC power) must still demonstrate through analysis that it meets cooling requirements under worst-case conditions. A molten salt-fueled reactor must explain how it meets safety criteria given a liquid fuel that doubles as coolant - a paradigm very different from LWRs. These differences can make the FSAR larger or more complex in certain areas, as the applicant might have to include additional background, experimental validation, or request exemptions from regulations that assumed LWR technology. NRC's regulatory architecture—particularly under Part 52—was built around conventional LWRs and may not fit novel technologies or modular deployment strategies without targeted reform  \cite{Righetti2025}. Indeed, one identified licensing barrier for SMRs is the prescriptive regulatory framework that was developed for conventional reactors and may impose requirements not fully applicable to new designs \cite{Sam2023}. In an FSAR, this often translates to the need to justify alternative approaches or criteria. For instance, an advanced reactor might request an exemption from certain emergency planning requirements by showing (via FSAR analysis) that its source term is so limited that offsite impact would be minimal.

\subsection{Risk-Informed and Performance-Based Approaches}
\vspace{0.2cm}

Both NRC and the international community (e.g., IAEA) are encouraging more risk-informed licensing, which can benefit Gen IV designs. Rather than purely following deterministic rules, risk-informed methods use probabilistic risk assessment (PRA) to focus on the most important safety issues. Advanced reactor FSARs may incorporate PRA insights to justify simpler safety systems or reduced redundancies, on the basis that certain severe accidents are rendered essentially impossible by design. The NRC has acknowledged this by adapting guidance like the “Licensing Modernization Project” framework for non-LWRs, which applicants can use to structure their safety analyses. As a result, an FSAR for a non-LWR reactor might have a different mix of analyses – for example, more on mechanistic source term calculations and PRA outcomes – compared to a traditional reactor FSAR. This is a tool, not a hindrance: it allows innovators to make a safety case in a way that suits their design, as long as they meet the overarching safety goals. The forthcoming 10 CFR Part 53 (a new framework for advanced reactors) explicitly aims to allow this flexibility, letting applicants organize their safety case around functional criteria and risk insights \cite{Righetti2025}. In practice, even before Part 53, NRC can use its existing regulatory tools (like custom license conditions or tailored review plans) to accommodate these novel approaches in the FSAR.

\subsection{SMR Advantages – Simpler FSAR?}
\vspace{0.2cm}

SMRs are often touted as simpler and inherently safer. Does that make the FSAR easier? In some respects, yes. For instance, NuScale's SMR design (a modular integral pressurized-water reactor) leverages passive cooling and a small core, which allowed for a smaller emergency planning zone and fewer complex safety systems. Their FSAR could demonstrate large safety margins in worst-case accidents, arguably making the NRC's evaluation more straightforward. Additionally, SMRs often plan to replicate modules, so once a design is certified or licensed, subsequent units should face minimal licensing hurdles. However, an FSAR for the first of a kind SMR still must be extremely detailed. NuScale's design certification application to the NRC comprised a 12,000-page FSAR (plus millions of pages of supporting documents), which took the NRC staff about 42 months of review to approve. The depth of analysis did not diminish – if anything, it included novel analyses (e.g., for multi-module interactions, and long-term cooling without operator action) that traditional reactors didn't require. The key point is that the FSAR framework was flexible enough to handle this: NuScale's FSAR made the case for its unique design features, and the NRC was able to conclude those features met safety requirements without having to fundamentally change the licensing process.

\subsection{IAEA and International Practices}
\vspace{0.2cm}

Globally, the fundamentals of safety analysis reports remain consistent. The IAEA's guidelines \cite{IAEA2021} encourage member states to adopt a graded approach: the effort and detail in an FSAR should be commensurate with the complexity and risk of the facility. For SMRs and advanced reactors, many countries are considering how to right-size their requirements. For example, Canada's regulator (CNSC) still requires a full safety case document, but is open to innovation in methodology. The United Kingdom’s Generic Design Assessment (GDA) process similarly involves a detailed Pre-Construction Safety Report, effectively an FSAR, for new designs like Rolls-Royce SMR or advanced reactors. These international processes are analogous to the NRC’s: they require comprehensive documentation of safety, but they also allow some tailoring. The IAEA plays a harmonizing role by publishing safety standards – like SSG-61 for SAR content and SSG-12 for licensing processes – which help ensure that an FSAR for an advanced reactor in one country will cover the same essential points as an FSAR in another. This comparability is useful not only for safety but also for efficiency: reactor vendors can reuse portions of their safety case across different jurisdictions. In other words, a well-prepared FSAR is a \textit{transferable asset}. Companies pursuing global deployment can leverage the work done for an NRC FSAR to satisfy regulators elsewhere, adjusting for local requirements but keeping the core safety arguments consistent.

\section{Case Studies: Lessons from Recent Licensing Efforts}

Experience from recent reactor licensing efforts underscores how critical the FSAR and the applicant's licensing strategy are to outcomes. Here we examine a few case studies, spanning both successes and setbacks, to draw out lessons on how the FSAR can be a tool to expedite licensing or, if handled poorly, a source of delay.

\subsection{NuScale SMR Design Certification (Success through Preparation)}
\vspace{0.2cm}

NuScale Power’s SMR is a 50 MWe (later uprated to ~60 MWe) light-water reactor module. NuScale was the first U.S. SMR developer to go through NRC’s full design certification process. They submitted their Design Certification Application (DCA) – essentially an FSAR for the standard design – in January 2017, and the NRC completed its final safety evaluation by September 2020, with the design formally certified in 2023 \cite{DOE2023}. This made NuScale’s design the first SMR cleared for use in the U.S., a milestone achievement. \\

NuScale’s FSAR had to demonstrate the safety of several novel features: natural circulation core cooling (no large primary pumps), a high degree of modularity (up to 12 modules at one plant), and passive safety systems that could keep the reactor stable without operator action for extended periods. The NRC issued hundreds of RAIs, but NuScale was well-prepared to answer them, having invested significant effort in testing and analysis. The company built a one-third scale test facility to generate data for the FSAR’s thermal-hydraulic models, and it worked closely with national labs and the DOE (with over \$400 million in DOE support) to resolve technical issues \cite{DOE2020}. The result was an application that, while large, was of high quality. The NRC noted that NuScale’s timely RAI responses and mature design helped keep the review on track \cite{NIA2024}. NuScale also smartly phased certain issues: for example, they deferred some design features to a later amendment (such as uprating the module power) so as not to delay the initial certification. This demonstrates strategic use of the licensing process as a tool – they got a baseline design approved, which can now be built, and will incorporate incremental improvements through license amendments. \\

In less than four years of review time, NuScale achieved NRC staff approval. The FSAR became the basis for the NRC’s safety evaluation report \cite{NUREG2246}, which confirmed that the design met all applicable safety requirements. The success of NuScale’s FSAR can be gauged by the fact that the Commission voted unanimously to certify the design – indicating confidence in the documentation and analysis provided. Going forward, any utility that applies for a Combined License to build a NuScale plant can reference the certified design’s FSAR, avoiding re-doing those analyses. This greatly streamlines the COL process for deployments. The NuScale case illustrates that while preparing an FSAR is an enormous undertaking (the application was tens of thousands of pages), doing it rigorously pays off. It turned the FSAR into a springboard for the company’s commercial plans, rather than a roadblock.

\subsection{Oklo Aurora Microreactor (Challenges of an Incomplete Application)}
\vspace{0.2cm}

Oklo, Inc. is a startup that in 2020 became the first private company to submit a Combined License Application for an advanced (non-light-water) microreactor. Their Aurora reactor is a very small (\~1.5 MWe) fast reactor with a heat pipe cooling system – a novel design with no direct precedents in NRC licensing. The application, including the FSAR, was accepted by NRC for review in June 2020. However, in January 2022 the NRC denied Oklo’s license application without prejudice (meaning Oklo could reapply) because of significant gaps in the safety information \cite{NRC2022Oklo}. The case is widely discussed as an example of how FSAR completeness determines whether a review can proceed. \\

In NRC’s denial letter and press release, the agency specified that "Oklo's application continues to contain significant information gaps in its description of Aurora’s potential accidents as well as its classification of safety systems and components" \cite{NRC2022Oklo}. In other words, two of the most fundamental parts of any FSAR – the accident analysis and the safety classification of hardware – were not adequately addressed. The NRC had engaged with Oklo for nearly 22 months, requesting additional information multiple times. Oklo provided some supplemental reports in 2021, but ultimately the NRC staff found these responses insufficient to close the gaps. For example, a reactor's FSAR must analyze the worst-case accidents (sometimes called Design Basis Accidents) and demonstrate how the reactor’s design either prevents these scenarios or mitigates their consequences. The NRC indicated Oklo had not convincingly analyzed certain possible accidents for their design. Additionally, every FSAR must categorize structures, systems, and components by their safety importance (safety-related vs. non-safety, etc., which affects their required reliability and quality assurance). Oklo's classification approach was apparently unclear or lacking, which is a serious deficiency because it underpins the whole safety argument.\\

Why did this happen? The application was the first of its kind: a custom combined license application for a non-light-water microreactor with no direct licensing precedent. It used a "maximum credible accident" approach that was not part of established NRC guidance. The NRC's denial identified gaps in two areas: the description of Aurora's potential accidents and the classification of safety systems and components \cite{NRC2022Oklo}. These are the parts of an FSAR on which the rest of the safety case depends. The NRC did try to work with them (the staff even took the unusual step of not issuing a review schedule, instead focusing on resolving key gaps early), but by late 2021 it became clear the necessary info was not forthcoming. Oklo ultimately did not contest the denial; instead, they went back to the drawing board to improve their design and application for a future re-submittal \cite{WNN2022}. \\

Oklo's initial application was stopped after about 2 years of effort, illustrating that an incomplete FSAR can indeed become a barrier – in this case halting the regulatory process entirely. The positive spin is that this happened early, before large amounts of NRC resources were spent on a futile review. The lesson for any first-of-a-kind applicant is that there is no shortcut around providing a full safety case. Any innovative reactor must still satisfy the regulator’s fundamental questions: “What accidents can happen, and how do you prevent or cope with them? What systems are crucial to safety, and how reliable are they?” These are the core of an FSAR. The Oklo case shows that if those questions are not fully answered in the application, the NRC cannot proceed. It underscores that the FSAR is the license application in many respects – if it is lacking, the process cannot move forward. On a constructive note, NRC’s door remains open to Oklo or others to reapply with a stronger FSAR. The incident also prodded discussions industry-wide about how NRC can better communicate expectations to first-time applicants, and whether guidance for novel microreactors needs to be clearer. It also highlighted the value of licensing experience across the industry: people who have written or reviewed FSARs can catch gaps of this kind before submission. In short, the case shows that the FSAR must be approached with the same engineering rigor as building the reactor itself.\footnote{Since this paper was written, Oklo has resumed pre-application work with the NRC. Oklo reported that the NRC's pre-application readiness assessment for Phase 1 of a new combined license application identified no significant gaps that would hinder acceptance \cite{NRC2025OkloReadiness, Oklo2025Readiness}.}

\subsection{Kairos Power Hermes Reactor (Iterative Success with a Graded Approach)}
\vspace{0.2cm}

Kairos Power, a venture-backed company, is developing a fluoride salt-cooled high-temperature reactor (KP-FHR), a type of Gen IV reactor. To incrementalize development, Kairos pursued a two-step approach. First, they applied for a Construction Permit for a low-power test reactor called Hermes (approx. 5 MWt) in 2021–2022. The NRC granted that construction permit by the end of 2023, making it the first non-LWR fission reactor to receive an NRC construction permit in decades. Building on that, Kairos immediately applied in mid-2023 for a larger test/demonstration reactor, Hermes 2 (two modules, each 35 MWt, capable of power generation). Impressively, the NRC issued the construction permits for Hermes 2 in November 2024 – just 16 months after the application, an accelerated timeline by any measure \cite{GAIN2024}. This case study shows how a well-planned licensing strategy and high-quality safety documentation can enable fast regulatory approval for advanced reactors. \\

Kairos invested heavily in pre-application engagement. Starting as early as 2018, they interacted with NRC through the agency’s optional "Risk-Informed Licensing Initiative" and submitted white papers on key topics (like their fuel and coolant). By the time they submitted the Hermes construction permit application (which included a Preliminary Safety Analysis Report, similar in content to an FSAR but for a test reactor), the NRC was already familiar with many aspects of the design. The initial Hermes CP was approved in roughly two years. For Hermes 2, the timeline was even shorter: application to final safety evaluation in about 1 year. Kairos achieved this by treating the first Hermes as a learning exercise both for themselves and the NRC. The Hermes 2 FSAR could then reference much of the data and analysis from Hermes 1, but scaled up and updated. The NRC noted that the extensive pre-application engagement and the fact that Hermes 2 was an evolution (not a brand-new concept) enabled an accelerated review. In essence, Kairos "de-risked" their FSAR by resolving many safety questions in stages. \\

Even as a test reactor, Hermes 2 had to address unique features: molten salt coolant, Tri-structural Isotropic \\(TRISO) particle fuel in pebble form, and a two-reactor configuration. Kairos’s FSAR broke these challenges down. For example, TRISO fuel has excellent fission product retention, which was a cornerstone of their safety case – the FSAR’s accident analyses could show very low offsite dose consequences even in severe events. The coolant is at atmospheric pressure (another safety plus compared to high-pressure water reactors), which simplified certain analyses. Kairos likely used these inherent advantages to argue that even if some systems failed, the reactor would remain within safe limits. The FSAR also classified safety systems in a risk-informed way, focusing on passive features. Because the NRC had been involved in methodology discussions earlier, there was mutual understanding that ended up minimizing RAIs. Kairos also made use of iterative hardware testing (they had operated non-nuclear prototypes of their coolant loops) to supply empirical evidence in the FSAR. All these choices reflect using the FSAR process as a development tool: start small, prove the safety step by step, document it thoroughly, and build regulator confidence. \\

Kairos's Hermes 2 construction permit approval is a success story for Gen IV reactor licensing. Achieving it in 16 months from application is nearly unprecedented for nuclear facilities. It shows that with the right approach, the FSAR and overall application for an advanced reactor need not take the better part of a decade; it can be done on a start-up timeline. Importantly, the NRC’s final safety evaluation for Hermes 2 concluded that the design, as presented, can be built safely. This gives Kairos the green light to construct and test their reactors, gathering data for an eventual operating license and larger commercial units. 

The case underscores a few general lessons: 
\vspace{0.2cm}

\begin{enumerate}
    \item \textbf{Engage early and often:} By the time the FSAR is submitted, no surprises should remain in the eyes of the regulator. Kairos’s years of pre-engagement exemplify this.
    \item \textbf{One bite at a time:} If your ultimate reactor is very novel, consider staging the licenses (e.g., test reactor then full-scale reactor) so that the first FSAR is manageable and builds trust. 
    \item \textbf{Leverage FSAR as a knowledge base:} Each licensing document can feed into the next. Kairos treated the FSAR as a living, upgradable tool – Hermes 1’s PSAR fed Hermes 2’s FSAR, which will likely feed into a full-scale plant FSAR down the line. This iterative approach can make the learning curve less steep for both applicant and regulator.
\end{enumerate} 

\section{Licensing Timelines, Delays, and the Influence of FSAR Quality}

One of the most common concerns for reactor developers is how long the licensing process will take. While there are many factors (including NRC staffing, public intervention, and policy changes), the quality of the FSAR and the expertise of the team preparing it stand out as factors within the applicant’s control that heavily influence timelines. \\

Legal scholars have argued that successful implementation of Part 53 will depend not only on rulemaking but also on internal cultural shifts at the NRC and clear guidance to industry \cite{Righetti2025}.

Historically, NRC reviews for new reactors have varied widely in length. For example, standard design certifications for large reactors in the 1990s (like the GE ABWR) took on the order of 5–6 years. The AP1000 design certification in the 2000s took about 3-4 years. In recent cases, we see NuScale's SMR reviewed in \~3.5 years (42 months) for the technical review phase, Oklo's application effectively stopped after \~1.5 years of review, and Kairos's test reactor permits approved in ~1–2 years \cite{GAIN2024}.\\ 

This suggests that smaller or well-prepared applications can indeed be reviewed faster, whereas incomplete or novel ones without sufficient support can face indefinite delays or rejection. FSAR Quality and RAIs: A strong correlation exists between the initial completeness of an FSAR and the number of RAIs or supplements needed. Each round of RAIs can add months as the applicant prepares responses and the NRC analyzes them. Industry experience indicates that timely RAI responses and minimizing the need for multiple rounds of questions are key to staying on schedule \cite{NIA2024}. For instance, NuScale's ability to respond to most RAIs within a 30-day window kept the NRC's review clock moving consistently. In contrast, if an FSAR has major gaps, the first RAIs might point to whole sections needing rewrite or new analysis, which can derail a schedule. Oklo's case epitomized this: the "RAIs" were essentially fundamental and could not be resolved in a timely way, leading to a termination of review. On the other hand, Kairos reportedly had relatively few RAIs on critical issues due to resolving them earlier, allowing the NRC staff to meet an aggressive schedule. 

\begin{itemize}
    \item \textbf{Team Expertise:} The composition of the applicant’s licensing team (and supporting contractors) often makes the difference in FSAR quality. Teams that include people who have navigated NRC licensing before (whether from industry or from retired NRC staff) can anticipate what information the NRC needs. They are more likely to write the FSAR in a manner that is clear and hits the regulatory marks. By contrast, newcomers might unknowingly omit needed details or rationale. \cite{Sam2023} identify "regulatory capability gaps" as a challenge for emerging SMR developers - essentially the learning curve of understanding regulatory expectations. Filling those gaps by hiring experienced consultants or partnering with larger organizations can mitigate the risk of an inadequate FSAR. This is one reason some newer reactor companies have partnered with DOE labs or even large vendors to assist in their licensing efforts. Moreover, NRC has been willing to have pre-application dialogues (for example, through its Vendor Design Review or topical report process) to guide new applicants – but it’s incumbent on the applicant to seek and utilize that guidance. 
    \item \textbf{Process Improvements and Future Outlook:} The NRC itself, along with industry, is seeking ways to make licensing more efficient without sacrificing rigor. One example is the development of the new Part 53 rule, which aims for a technology-inclusive, risk-informed framework by 2025. While Part 53 won't eliminate the need for an FSAR-equivalent safety case, it might allow a more streamlined structure (for instance, organizing the FSAR by safety functions rather than prescriptive systems). The hope is that this could reduce unnecessary information submittals and focus documentation on what truly matters for safety. Additionally, NRC has experimented with "standard design approvals" and other mechanisms to give partial approvals early. None of these changes remove the fundamental task of demonstrating safety in writing – they simply offer different pathways to present and judge that information. The IAEA’s Nuclear Harmonization and Standardization Initiative (launched in 2022) is also looking to align regulatory requirements internationally, which could mean a future where a single FSAR could serve in multiple countries with minimal modification. If successful, this would greatly amplify the value of investing in a high-quality FSAR. 
    \item \textbf{FSAR as a Management Tool:} Some companies have reframed how they view the FSAR internally. Instead of seeing it as a one-time document prepared by a siloed "licensing group," they treat it as a central project deliverable that integrates inputs from design, engineering, R\&D, and operations teams. This approach ensures that as the design evolves, the safety basis is evaluated in tandem. Modern software tools and configuration management can help keep the voluminous data organized. By the time of submission, such an FSAR is not a retroactive documentation of the design – it is a real-time reflection of it. This integration can prevent painful revisions and surprises. It also means the organization truly understands its own design’s risks and mitigations, rather than discovering them via NRC questioning. In effect, the process of writing the FSAR becomes a form of internal quality assurance and risk management.
\end{itemize}

\section{Improving Licensing Efficiency through NEPA Process Reform}

The licensing of nuclear reactors, particularly SMRs and advanced Generation IV reactors, can face significant delays due to lengthy environmental review processes required by the National Environmental Policy Act (NEPA). Recent efforts have focused on reforming these NEPA processes to reduce unnecessary regulatory burdens while ensuring environmental protection, thus expediting reactor deployment.

\subsection{Categorical Exclusions and Streamlined NEPA Reviews}

One prominent strategy to enhance licensing efficiency has been the consideration of categorical exclusions (CEs) for certain classes of reactors, particularly microreactors. Idaho National Laboratory (INL) has identified opportunities for using categorical exclusions under NEPA specifically tailored to microreactors, which by definition have a significantly lower environmental footprint compared to traditional nuclear power plants. These exclusions acknowledge that certain reactor types inherently involve fewer environmental impacts, and therefore detailed Environmental Impact Statements (EIS) or Environmental Assessments (EA) may be unnecessary. Utilizing CEs for microreactors can dramatically shorten the NEPA review timelines, enabling faster licensing and deployment without compromising environmental safeguards. INL's assessment emphasizes that clear criteria must be established to ensure that any microreactor eligible for categorical exclusion indeed has negligible impacts. These criteria include limited site disturbance, minimal emissions, and significantly reduced water usage \cite{inl2025categorical}. Establishing such categorical exclusions could lead to substantial reductions in both licensing duration and associated costs, allowing developers to allocate resources more efficiently to safety and technology improvements.

\subsection{Recent Improvements to the NRC's Mandatory Hearing Process}

Another crucial area of reform has been the streamlining of the NRC mandatory hearing process. Historically, the mandatory hearing has been a lengthy and resource-intensive step, requiring the Commission to independently confirm the findings of NRC staff before licensing approval. Recent NRC policy updates have significantly refined this process, clarifying and narrowing the scope of these hearings. Specifically, the NRC revised the mandatory hearing guidelines to reduce redundancy by allowing the Commission to rely more substantially on detailed staff analyses. The updated process emphasizes the use of existing documentation, such as Safety Evaluation Reports and Environmental Impact Statements, instead of generating new analyses specifically for hearings. This approach reduces repetitive work and expedites decision-making without sacrificing rigorous safety and environmental reviews. Moreover, the reform provides clearer guidelines for commissioners on focusing questions during hearings, prioritizing significant safety and environmental concerns, and limiting peripheral discussions. These adjustments have been projected to considerably reduce both the duration and complexity of the mandatory hearing process, thereby accelerating overall reactor licensing timelines.

\subsection{Implications for Advanced Reactor Licensing}

These NEPA reforms collectively hold great promise for enhancing the efficiency of advanced reactor licensing. By adopting categorical exclusions for microreactors and refining the mandatory hearing process, the NRC is positioning itself to handle the anticipated influx of new reactor applications more effectively. These steps not only streamline regulatory processes but also align the U.S. nuclear regulatory framework more closely with international best practices, facilitating global competitiveness and innovation in advanced reactor technologies.

\section{Conclusion}

The Final Safety Analysis Report (FSAR) remains one of the most critical deliverables in nuclear licensing, not simply as a regulatory requirement, but as a central engineering and safety document that defines a reactor's safety case. This paper has shown that the FSAR, when approached proactively, serves as a tool for collaboration, design validation, and long-term safety assurance. Yet in practice, many in the nuclear industry still treat it as a bureaucratic formality—an attitude mirrored across other safety-critical industries when documentation is viewed as compliance overhead rather than a cornerstone of quality. Just as the automotive sector eventually came to embrace the PPAP process, the nuclear sector must reframe its relationship with the FSAR.\\

To succeed in licensing next-generation technologies like SMRs and Gen IV reactors, the nuclear industry must do more than wait for regulatory reform. Although recent advancements—such as proposed NEPA streamlining, NRC hearing reforms, and the development of 10 CFR Part 53—offer long-term promise, these policy changes will take years to fully implement. What can change immediately is mindset. A cultural shift toward viewing documentation as a shared language of safety, not a chore, can begin now, and its effects can be felt long before new regulations are codified.\\

This cultural transformation, however, hinges on people. The industry must prioritize building a skilled pipeline of professionals capable of producing and maintaining rigorous FSARs and UFSARs. This is not glamorous work, and its long training times, combined with a lack of recognition and institutional support, make it an unattractive career path to many. Yet it is precisely in this area, engineering rigor, safety analysis, document management where the sector must invest. Without qualified personnel, even the best-intentioned reforms or innovations risk failure. Programs like DOE’s workforce development initiatives and IAEA’s Systematic Approach to Training offer templates, but the onus is on developers and utilities to make these roles visible, respected, and rewarding.\\

The FSAR, ultimately, is a reflection of an organization’s culture. A clear, complete, and thoughtful FSAR signals not only technical readiness but institutional maturity. When teams engage with the process early, respond thoroughly to RAIs, and use documentation as a development tool—not just a licensing artifact—they demonstrate a mindset of safety and accountability. This approach pays dividends in licensing speed, regulatory trust, and long-term operational excellence.\\

The message to advanced reactor developers is clear:

\begin{itemize}
    \item Do not defer your licensing strategy—embed it into design from day one.
    \item Treat the FSAR as a living design tool, not a compliance checkbox.
    \item Empower your people—build teams that understand not just physics but regulation, systems, and the art of good writing.
    \item Push for reform, but don’t wait for it—change begins with culture.
\end{itemize}

In an era demanding both speed and rigor, the FSAR can either slow you down or become your most powerful asset. Reframe it. Invest in it. Let it speak not just for your design, but for your readiness to bring safe, innovative nuclear energy to the world.

\bibliography{bibliographie.bib}
%\printbibliography

\end{document}


